\section{作为生成函数的作用量泛函}
\label{sec:action-functional-generating-function}

我们先把第一变分公式 (3.4.15) 限制到余切丛的情形。此时 $\alpha=-\theta$，而\personLiouville{} 1-形式为 $\theta=\sum_{i=1}^n p_i\,dq^i$，于是
\begin{equation}
\begin{split}
  d\mathcal A_H(\gamma)\xi
  ={}&\int_0^1\omega_0\bigl(\dot\gamma(t)-X_H(t,\gamma(t)),\xi(t)\bigr)\,dt
      +\bigl\langle\theta(\gamma(1)),\xi(1)\bigr\rangle \\
     &-\bigl\langle\theta(\gamma(0)),\xi(0)\bigr\rangle .
\end{split}
\label{eq:cotangent-action-first-variation}
\end{equation}
这里 $\gamma:[0,1]\to T^*N$，而 $\xi\in\Gamma(\gamma^*T(T^*N))$ 是沿 $\gamma$ 的一个向量场。

对这个公式稍加思考，就会得到若干重要结论。

首先，我们考虑始于\termZeroSection{}中某一点的道路 $\gamma:[0,1]\to T^*N$。记
\[
  \Omega(0;o_N)=\{\gamma\mid\gamma(0)\in o_N\}.
\]
有一个自然映射
\[
  \pi\operatorname{ev}_1:\Omega(0;o_N)\longrightarrow N,
\]
它由
\begin{equation}
  \pi\circ\operatorname{ev}_1(\gamma)=\pi(\gamma(1))
  \label{eq:projection-evaluation}
\end{equation}
定义。这给出了 $\Omega(0;o_N)$ 在 $N$ 上的一个纤维化，它在 $q\in N$ 上的纤维为
\[
  \Omega(o_N,T_q^*N)
  =\{\gamma\mid\gamma(0)\in o_N,\ \gamma(1)\in T_q^*N\}.
\]
\stfhIndex{I078}{Xianweidaoshu}{纤维导数（fiber derivative）}
$\mathcal A_H$ 相对于 $\pi\operatorname{ev}_1$ 在 $q$ 处的纤维导数，无非是
\[
  \mathcal A_H:\Omega(o_N,T_q^*N)\longrightarrow\mathbb R
\]
的第一变分。因而
\[
  d^{\mathrm{fiber}}\mathcal A_H(\gamma)(\xi)
  =\int_0^1\omega_0\bigl(\dot\gamma(t)-X_H(t,\gamma(t)),\xi(t)\bigr)\,dt
\]
\stfhIndex{I019}{00-sigma-ah}{$\Sigma_{\mathcal A_H}$}
\stfhIndex{I077}{Xianweilinjieji}{纤维临界集（fiber critical set）}
对所有 $\xi\in T_\gamma\Omega(o_N,T_q^*N)$ 成立。因此，相应的\termFiberCriticalSet{}（记为 $\Sigma_{\mathcal A_H}\subset\Omega(0;o_N)$）由\personHamilton{}方程的解组成：
\[
  \dot x=X_H(t,x),\qquad x(0)\in o_N,\quad x(1)\in T_q^*N.
\]
由式~\eqref{eq:cotangent-action-first-variation}，注意在 $\Sigma_{\mathcal A_H}$ 上有
\[
  d\mathcal A_H(\gamma)(\xi)
  =\bigl\langle\gamma(1),d\pi(\xi(1))\bigr\rangle.
\]

\begin{exercise}
\label{ex:push-forward-action-one-form}
补全一个形式的启发式论证，由这个公式推出：1-形式 $d\mathcal A_H(\gamma)$ 的\termPushForward{}无非是 $\gamma(1)\in T_q^*N$。
\end{exercise}

\stfhIndex{I197}{Wensitan}{温斯坦（Weinstein）}
这就完成了下述命题的启发式证明。该命题是\personWeinstein{}观察到的。

\begin{proposition}[\personWeinstein{}]
\label{prop:action-functional-generates-image-zero-section}
在上述意义下，二元组 $(\mathcal A_H,\Omega(0;o_N))$ 是 $\phi_H^1(o_N)$ 的一个生成函数。
\end{proposition}

事实上，这个生成函数有一个自然的\termFiniteDimensionalReduction{}。我们称它为 $L_H=\phi_H^1(o_N)$ 的\termBasicGeneratingFunction{}，并记为 $h_H:L_H\to\mathbb R$。对给定的 $x\in L_H$，记
\[
  z_x^H(t)=\phi_H^t\bigl((\phi_H^1)^{-1}(x)\bigr).
\]
它是一条\termHamiltonianTrajectory{}，且依定义满足
\begin{equation}
  z_x^H(0)\in o_N,\qquad z_x^H(1)=x.
  \label{eq:basic-hamiltonian-trajectory-endpoints}
\end{equation}
我们回顾 (Oh97b) 中用到的如下基本引理，其证明留作习题。

\begin{lemma}
\label{lem:basic-generating-function}
定义函数 $h_H:L_H\to\mathbb R$ 为
\[
  h_H(x)=\mathcal A_H(z_x^H),
  \qquad
  z_x^H(t)=\phi_H^t\bigl((\phi_H^1)^{-1}(x)\bigr).
\]
则它满足
\[
  i_H^*\theta=dh_H.
\]
也就是说，$h_H$ 是 $L_H$ 的一个生成函数。我们称 $h_H$ 为 $L_H$ 的\termBasicGeneratingFunction{}。
\end{lemma}

\begin{exercise}
\label{ex:prove-basic-generating-function}
证明这个引理。
\end{exercise}

公式~\eqref{eq:cotangent-action-first-variation} 的另一个结论，是它刻画了 $\mathcal A_H$ 的变分理论中的\termNaturalBoundaryCondition{}。一个显然的\termNaturalBoundaryCondition{}是周期边界条件 $\gamma(0)=\gamma(1)$。下面给出一类更有趣的\termNaturalBoundaryCondition{}。

\begin{proposition}
\label{prop:tautological-form-vanishes-on-conormal}
对子流形 $S\subset N$，考虑余法丛 $i_S:\nu^*S\hookrightarrow T^*N$。则 $i_S^*\theta=0$。
\end{proposition}

\begin{proof}
取 $\xi\in T_\alpha(\nu^*S)$，并设 $\pi(\alpha)=x$。\footnote{译者注：原书此处写作 $\xi\in T_\alpha(T^*N)$，但依 $i_S^*\theta$ 的定义及下一句，应为 $\xi\in T_\alpha(\nu^*S)$；紧随其后的公式还脱漏了最外层右括号。译文据此补正。}我们有
\[
  i_S^*\theta(\xi)=\alpha(x)\bigl(d\pi(\xi(\alpha))\bigr).
\]
但这个配对为零，因为依定义有 $d\pi(\xi(\alpha))\in T_xS$ 和 $\alpha(x)\in\nu_x^*S$。证明完毕。
\end{proof}

由该命题，若把作用量泛函 $\mathcal A_H$ 限制到子集
\[
  \Omega_{S_0S_1}
  =\Omega(\nu^*S_0,\nu^*S_1)
  =\{\gamma:[0,1]\to T^*N
      \mid\gamma(0)\in\nu^*S_0,\ \gamma(1)\in\nu^*S_1\},
\]
则 $\mathcal A_H$ 的第一变分 $d\mathcal A_H$ 化为
\begin{equation}
  d\mathcal A_H(\gamma)(\xi)
  =\int_0^1\omega_0\bigl(\dot\gamma(t)-X_H(t,\gamma(t)),\xi(t)\bigr)\,dt.
  \label{eq:conormal-action-first-variation}
\end{equation}
我们指出，上面定义的道路空间 $\Omega(o_N,T_q^*N)$ 是余法边界条件的一个特例，它对应于 $S_0=N$ 和 $S_1=\{q\}$。

\cref{prop:tautological-form-vanishes-on-conormal}的一个立即推论是对临界点方程的如下刻画：$d\mathcal A_H|_{\Omega(S_0,S_1)}(\gamma)=0$ 当且仅当 $\gamma$ 满足
\begin{equation}
  \dot\gamma=X_H(t,\gamma(t)),
  \qquad
  \gamma(0)\in\nu^*S_0,\quad \gamma(1)\in\nu^*S_1.
  \label{eq:conormal-hamilton-equation}
\end{equation}

\begin{definition}[作用谱]
\label{def:conormal-action-spectrum}
我们把 $\operatorname{Spec}(H;S_0,S_1)$ 定义为临界值的集合，并称它为 $\mathcal A_H$ 在 $\Omega_{S_0S_1}$ 上的\termActionSpectrum{}，即
\[
  \operatorname{Spec}(H;S_0,S_1)
  =\{\mathcal A_H(z)\mid
      \dot z=X_H(t,z(t)),\ z(0)\in\nu^*S_0,\ z(1)\in\nu^*S_1\}.
\]
\end{definition}

\begin{proposition}
\label{prop:conormal-action-spectrum-measure-zero}
子集 $\operatorname{Spec}(H;S_0,S_1)\subset\mathbb R$ 是紧的，且测度为零。
\end{proposition}

\begin{proof}
如第 3.1 节所示，式~\eqref{eq:conormal-hamilton-equation} 的解与交集
\[
  \phi_H^1(\nu^*S_0)\cap\nu^*S_1
\]
之间存在一一对应。显然，这个集合是紧的。\footnote{译者注：原书在没有附加横截性或紧性假设时断言该交集“显然”紧；例如取 $H=0$ 且 $S_0=S_1=\{q\}$，该交集就是非紧的余切纤维 $T_q^*N$。因此原书此处的紧性论证不能按所写成立。}另一方面，从 $\nu^*S_1$ 到 $\mathbb R$ 的函数
\[
  h:x\longmapsto z_x^H\longmapsto\mathcal A_H(z_x^H)
\]
是光滑的，因而它在 $\phi_H^1(\nu^*S_0)\cap\nu^*S_1$ 上的像是紧的。\footnote{译者注：此前的 $z_x^H$ 与 $h_H$ 是针对 $x\in\phi_H^1(o_N)$ 定义的；此处却在一般 $S_0$ 下把 $h$ 写成 $\nu^*S_1$ 上的函数，并在下述习题中把其临界点集等同于 $\phi_H^1(\nu^*S_0)\cap\nu^*S_1$。原书没有给出完成这一推广所需的新定义或论证。}

\begin{exercise}
\label{ex:critical-points-basic-function}
$h$ 的临界点集与交集 $\phi_H^1(\nu^*S_0)\cap\nu^*S_1$ 重合，而 $h$ 的临界值集与 $\operatorname{Spec}(H;S_0,S_1)$ 重合。
\end{exercise}

把\personSard{}定理应用于光滑函数 $h:\nu^*S_1\to\mathbb R$，即完成证明。
\end{proof}

我们还有如下\termTransversality{}结果：
\[
  \phi_H^1(\nu^*S_0)\pitchfork\nu^*S_1
\]
当且仅当式~\eqref{eq:conormal-hamilton-equation} 的\termLinearizationOperator{}
\begin{equation}
  \nabla_{\gamma'}-DX_H(\gamma):
  \Gamma_{S_0S_1}\bigl(\gamma^*(T(T^*N))\bigr)
  \longrightarrow
  \Gamma\bigl(\gamma^*T(T^*N)\bigr)
  \label{eq:conormal-linearization}
\end{equation}
是满射。这里 $\Gamma_{S_0S_1}(\gamma^*(T(T^*N)))$ 是 $\Gamma(\gamma^*T(T^*N))$ 的子集，定义为
\[
\begin{split}
  \Gamma_{S_0S_1}\bigl(\gamma^*(T(T^*N))\bigr)
  =\{\xi\in\Gamma\bigl(\gamma^*T(T^*N)\bigr)\mid{}&
    \xi(0)\in T_{\gamma(0)}(\nu^*S_0),\\
   &\xi(1)\in T_{\gamma(1)}(\nu^*S_1)\}.
\end{split}
\]
子集 $\Gamma_{S_0S_1}(\gamma^*(T(T^*N)))$ 也是 $\Omega_{S_0S_1}$ 在 $\gamma$ 处的切空间 $T_\gamma\Omega_{S_0S_1}$。这一对应是联系\personHamilton{}方程的\emph{动力学}与拉格朗日交\emph{几何}的关键。
